\section{Design Goal Formulation}
\label{sec:design_goal}
We now present a contextual inquiry of data science workflows within the HRKG platform 
%at \company. We discuss their workflow challenges and limitations while situating these observations in the context of existing studies of data science workflows. The study 
that helped understand the challenges of the practitioners with their existing workflows and informed design guidelines for building assistive tools.

\subsection{Contextual Inquiry of Graph Exploration Workflows}
Contextual inquiries are ethnographic field studies where researchers interview a small group of people in their natural environment as they conduct their work activities (\textit{context}) to gain an in-depth understanding of their work practices~\cite{holtzblatt1997contextual}. 
%The study method is particularly suitable for the industry setting, focusing on understanding expert users’ interactions with complex systems, the in-depth processes involved, and their mental models. 
As opposed to formative studies which are reflective and require participants to recall their experiences from memory, contextual inquiries take place in the participants' day-to-day work environment --- \company in this case.
% \hkc{suggestion: divide the subsection into two paragraphs:  why contextual inquiry (above) -> details of the study (below). add who was the subjects and why they were chosen and describe session}
We conducted the study by recruiting three data practitioners --- one researcher, one engineer, and one data scientist --- by reaching out to them in Slack\footnote{\url{https://slack.com/}}. 
These participants contributed to various projects related to knowledge acquisition (discovering new entities and relations), curation (graph understanding and standardization), and integration (graph alignment and merging) within the HRKG platform. Participants in each session brought their laptops, tools, and documents (Google Sheets and Slides).
The sessions lasted for $45$ minutes, where we observed participants performing their project-specific tasks with their existing tools and setup.
% and intervened as needed to ask follow-up questions.    
During observation, we intervened as needed to ask follow-up questions to gain a deeper understanding of the participants'
decisions and their pain points. 

% We further explained our interpretations of their tasks for clarification purposes. Before presenting the study observations, we briefly discuss the work setup at \company with a motivating example.

%\stitle{Graph construction and serving.}

% \todo{Contextual inquiry~\cite{holtzblatt1997contextual} is a type of ethnographic field study that involves in-depth observation and interviews of a small sample of users to gain a robust understanding of work practices and behaviors. Its name describes exactly what makes it valuable — inquiry in context:
% Context: The research takes place in the users’ natural environment as they conduct their activities the way they normally would. The context could be in their home, office, or somewhere else entirely.
% Inquiry: The researcher watches the user as she performs her task and asks for information to understand how and why users do what they do.
% Contextual inquiry is useful for many domains, but it is especially well-suited for understanding users’ interactions with complex systems and in-depth processes, as well as the point of view of expert users.}

\subsection{Key Takeaways}
\stitle{Tedious Context Switching.} While the participants were involved in different projects within the platform, such as knowledge acquisition, curation, and integration, their work setup involving graphs was somewhat similar. These individual projects or phases could be further broken down into smaller tasks ---  viewing and assessing data, formulating a hypothesis on how to perform the intended task, and verifying the outcomes of the executed task --- similar to Rahman et al.~\cite{rahman2022ie}. However, completing such multi-purpose tasks required participants to use different tools --- command-line environments, computational notebooks such as Jupyter, and code editors such as Visual Studio Code\footnote{https://code.visualstudio.com/} for programming; Neo4j browser\footnote{\url{https://neo4j.com/developer/neo4j-browser/}} for querying and visually exploring graphs and tools such as Matplotlib~\cite{hunter2007matplotlib} for visualizing data;  editors such as spreadsheets, and notebook cell outputs for viewing and verifying program outputs. Participants complained about tedious context switching among these tools. 

\stitle{Lack of Flexibility in Workflows.} Surprisingly, participants' graph exploration experiences were limited to using the Neo4j browser. While there are bespoke tools~\cite{graphistry, graphileon,linkurious} and even Neo4j desktop extensions~\cite{neobloom, neodash,yworks}, using those meant (a) ``\emph{adding another new tool}'' to the already iterative and messy experience and (b) learning a new platform. Moreover, using desktop-based tools had additional limitations, \eg company policy prohibiting the hosting of proprietary and sensitive data on personal computers. 
When using the Neo4j browser the participants could not leverage rich graph layouts, visualizations, and interactions available in various charting libraries mentioned in Section~\ref{sec:related}. Additionally, the browser lacked support for common visualizations such as bar charts and scatter plots, which are essential for downstream analysis and decision making. For viewing these charts, participants used a combination of notebooks or code editors and charting libraries in Python.

% \stitle{Reusability and reproducibility.}
% While using the Neo4j browser for graph querying, participants mentioned the challenge of reusability. One participant copied the queries issued in the browser in a notepad for future use, which was cumbersome. The alternative solution --- \ie writing these queries in Jupyter notebook using libraries such as py2neo\footnote{\url{https://py2neo.org/2021.1/}} and saving the session --- hindered exploration due to a lack of robust interactive graph exploration libraries in notebooks. This observation echoes the takeaways from our survey of graph exploration literature in Section~\ref{sec:related_graph}. Similar to the observation of Batch and Elmqvist ~\cite{batch2017interactive}, one participant used their explorations in Neo4j browsers as a mechanism to capture and persist observations. Neo4j browser visualizations are essentially graph query results and are downloadable as either images or files (\eg CSV and JSON). However, the participant mentioned that the process was cumbersome. Moreover, the onus was on the participant to manually compile the provenance of the exploration session, had they wanted to.

% \stitle{Lack of advanced customization.}
% Another limitation related to only accessing information through query results  --- either in the browser or in notebooks --- was that the participants could not export specific observations within the query result or corresponding visualizations. One participant mentioned this was a ``typical experience'' that required writing post-processing functions in notebooks or scripts to extract the desired value. This action is akin to customizing the default behavior (\ie output) of a data operation (\ie query in this case). The other alternative was to refine the query parameters issued in the browser to filter the desired value. Participants of other data science workflow studies also requested such customizations~\cite{rahman2022ie, teddy2020chi}.  
 
%  then introduce graph as an important workflow and say how we use it. provide a hypothetical graph management workflow by drawing parallel to Saga at Apple~\cite{ilyas2022saga} and IBM content service system~\cite{bharadwaj2017creation}, Knowledge Vault at Google~\cite{knowledgevault}. metaphactory~\cite{haase2019metaphactory} 

% \subsection{Existing Solutions and Improvement Opportunities}
% While recent solutions~\cite{wu2020b2, bauerle2022symphony, kery2020mage, rahman2020leam} address many of the challenges outlined in Section~\ref{sec:data_science_limitation}, majority of these systems exhibit several gaps (see Table~\ref{tab:summary}). We now discuss the potential improvement opportunities with these tools which inform our design goals.

% \begin{table}[!htb]%\scriptsize
% %\vspace{-10pt}
% \begin{tabular}{c|ccccc}
% \textbf{Needs} & \textbf{B2} & \textbf{Mage} & \textbf{Symphony} & \textbf{Leam} & \textbf{GUESS}\\ \hline
% Knowledge Graph Visualization &  x  &  x  &  x  & x & \checkmark\\ \hline\\  
% Code and visualization blending &  \checkmark  &  \checkmark  &  \checkmark  & \checkmark & \checkmark\\ \hline\\
% Bespoke solution &  x  &  x  &  x  & \checkmark & \checkmark\\ \hline\\
% State reusability &  --  &  --  &  \checkmark  & -- & x\\ \hline\\
% Interaction provenance &  --  &  --  &  x  & -- & --\\ \hline\\
% Interactions Customizing Code Parameters &  x  &  --  &  x  & x & x\\
% \end{tabular}
% \caption{Limitations of existing tools.}
% \label{tab:summary}
% \vspace{-10pt}
% \end{table}

\subsection{Design Goals}
We opted for developing \system, an interactive widget library designed to specifically support the knowledge curation and integration tasks within computation notebooks. Through \system we aim to support the rendering of graph visualizations in notebooks and provide libraries that enable users to execute a suite of graph-specific data operations for visualizing and exploring graphs. \system is built on top of the \basesystem ~\cite{choi2023towards}, an extensible framework which enables data science workflows within interactive programming environments. Our choice was informed by the observations of the contextual inquiry and solidified by the recommendation by existing work to blend the expressivity of codes with the interactivity of visualizations to ensure flexibility in workflows~\cite{futzing2019moseying} and evaluation studies conducted in recent work highlighting the suitability of such environments in eliminating context switching~\cite{wu2020b2, bauerle2022symphony}. 

%\stitle{D1. Enable in-notebook graph exploration.} 


% \stitle{D2. Provide support for interaction provenance.} As user interactions trigger data operations that update the visualizations, \system should track both fine-grained interaction details and the corresponding data state of the visualization re-rendered due to an update. Moreover, the interaction history, \ie provenance, should be easily accessible and tightly coupled with the visualization. 

% \stitle{D3. Ensure reusability of data/visualization states.} \system should provide mechanisms to access any data or visualization states within the user's current analytics session so that the state can be used for subsequent steps within or across sessions and projects. 

% \stitle{D4. Enable on-demand customization of data operations.}
% \system should provide flexible affordances to customize the underlying data operations of a visual interaction so that users can tune their exploration objectives and reflect those on the corresponding interface.